% TOP_PROBLEM: {"id":30007550,"problem_number":"LOCAL-30007550","title":"Demography and long-run growth","category_id":21,"category":{"id":21,"name":"economics","display_name":"Economics","description":"Open economic theory statements, published research agendas, and proposed scoped specifications; question provenance and historical priority are qualified per record.","slug":"economics","order_index":21,"created_at":"2026-10-08T00:00:00Z"},"status":"provenance_pending","external_url":null,"legacy_ids":[],"published":false,"statement":"Identify the research elasticities and effective-effort assumptions needed for credible growth predictions under declining population; the standard fixed-parameter growth identity is already known.","economics":{"original_id":39,"field":"Growth and productivity","kind":"identification","formalization_basis":"proposed_specification","source_status":"not_verified","priority_status":"not_established","priority_note":"No explicit primary open-problem statement matching this catalogue item was verified; supporting research is not evidence of first formulation.","earliest_verified_reference":null,"current_reference":{"authors":["Charles I. Jones"],"title":"The Past and Future of Economic Growth: A Semi-Endogenous Perspective","year":2022,"url":"https://www.annualreviews.org/content/journals/10.1146/annurev-economics-080521-012458","doi":"10.1146/annurev-economics-080521-012458","locator":"Abstract, especially final paragraph on untapped researchers and AI augmentation","evidence_summary":"The review gives the semi-endogenous relation between research-effort growth and long-run growth and names possible offsets from talent and AI. For specified parameters the basic demographic model is answered; the broader empirical generalisation is not a canonical unsolved theorem."},"formalization_note":"The broad catalogue item is converted into an empirical generalization problem. The source already explains the standard model; no exact open statement of this parameter-identification target was verified."},"statusQualification":"Provenance pending: an explicit published statement of this exact open question was not verified. This is a catalogue-authored candidate specification, not a certified open conjecture. The broad catalogue item is converted into an empirical generalization problem. The source already explains the standard model; no exact open statement of this parameter-identification target was verified.","statusReviewedAt":"2026-10-08"}
% FIRST_POSTED: 2026-10-08
% ENTRY_KIND: statement_only
% !TeX program = lualatex
\documentclass[11pt]{article}
\usepackage{fontspec}
\usepackage[a4paper,margin=25mm]{geometry}
\usepackage{amsmath,amssymb,mathtools}
\usepackage{xurl}
\usepackage[hidelinks]{hyperref}
\newcommand{\sref}[2]{\hyperlink{source:#1:#2}{[S#2]}}
\setlength{\emergencystretch}{3em}
\begin{document}
% problemId:30007550
\section*{Demography and long-run growth}\label{30007550}
\subsection{Definitions and mathematical statement}
Consider the semi-endogenous idea-production relation
\[
\dot A_t=\chi R_t^\lambda A_t^\phi,\qquad
R_t=N_t s_t h_t z_t,
\]
where \(A_t>0\) is useful knowledge, \(R_t\) effective research effort, \(\chi>0\), \(\lambda>0\), and \(\phi<1\). Population \(N_t\), researcher share \(s_t\), research skill \(h_t\), and technological augmentation \(z_t\) determine effort; output per person uses a stated production relation with \(A_t\). For fixed parameters and positive balanced effort growth \(g_R>0\), when a positive balanced-growth path exists, the model’s familiar implication is \(g_A=\lambda g_R/(1-\phi)\); establishing that algebraic relation is not the unresolved target.
Identify a credible joint set for \((\lambda,\phi)\) and future effort components under an explicitly declared declining-population scenario. Distinguish research quality, endogenous allocation into research, and AI augmentation from raw researcher headcounts. Use independent innovation measures and variation in research supply to assess whether historically estimated elasticities transport to demographic contraction. Compute corresponding bounds for long-run per-capita growth or for finite transition paths when balanced growth does not exist. State sensitivity to the production relation and to knowledge obsolescence. The target is empirical identification and scenario validity; demographics alone do not determine a model-free growth forecast, and known special-case demographic growth results must not be mislabeled unsolved.

\par\textbf{Formulation and scope.} The broad catalogue item is converted into an empirical generalization problem. The source already explains the standard model; no exact open statement of this parameter-identification target was verified.

\par\textbf{Open-question provenance.} An explicit published open-question statement was not verified. Sources below are supporting literature and must not be treated as a first listing.

\par\textbf{Historical priority.} No explicit primary open-problem statement matching this catalogue item was verified; supporting research is not evidence of first formulation.

\subsection{Short English statement}
Identify the research elasticities and effective-effort assumptions needed for credible growth predictions under declining population; the standard fixed-parameter growth identity is already known.

\subsection{Sources}
\begin{itemize}\raggedright
\item[\textnormal{[S1]}]\hypertarget{source:30007550:1}{} Charles I. Jones. 2022. The Past and Future of Economic Growth: A Semi-Endogenous Perspective. Abstract, especially final paragraph on untapped researchers and AI augmentation.
\url{https://www.annualreviews.org/content/journals/10.1146/annurev-economics-080521-012458}.
DOI: 10.1146/annurev-economics-080521-012458.
{\small\textit{Supporting or current literature; see evidence qualification. The review gives the semi-endogenous relation between research-effort growth and long-run growth and names possible offsets from talent and AI. For specified parameters the basic demographic model is answered; the broader empirical generalisation is not a canonical unsolved theorem.}}
\end{itemize}
\end{document}
